%%%%%%%%%%%%%%%%%%%%%%%%%%%%%%%%%%%%%%%%%%%%%%%%%%%%%%%%%%%%%%%%%%%%%%%%%%%%%%%%
% Diese Datei beinhaltet den eigentlichen Inhalt Ihrer Arbeit.
%
% Es bietet sich der Übersicht halber an, die einzelnen Abschnitte jeweils
% in eigene Dateien zu schreiben und mittels \input einzubinden.
% Eine mögliche Verzeichnisstruktur sähe entsprechend so aus:
%
%     thesis/
%     +- tex/
%     |  +- introduction.tex
%     |  +- motivation.tex
%     |  +- experiments.tex
%     |  |  ...
%     |  +- conclusion.tex
%     +- abstract.tex
%     +- contents.tex
%     +- thesis.tex
%%%%%%%%%%%%%%%%%%%%%%%%%%%%%%%%%%%%%%%%%%%%%%%%%%%%%%%%%%%%%%%%%%%%%%%%%%%%%%%%

\section{Einleitung}

\subsection{Motivation}
Im Modul \emph{ProgrammierPraktikum 2} wurde Test-Driven Development (TDD) eingeführt. Dabei erhielt ich erstmals Einblicke in praktische Teststrategien, insbesondere in die \emph{Testpyramide}. Als Anfänger faszinierten mich die klare Struktur und die methodische Vorgehensweise. Dies veranlasste mich dazu, die Frage zu stellen, ob dieser Ansatz universell angewendet wird oder ob alternative Modelle existieren.

Aus meinem aktuellen Werkstudentenjob ergab sich eine weitere Beobachtung. In einigen Projekten werden vorwiegend Systemtests implementiert, die das Gesamtverhalten der Anwendung abdecken. Separate Unit-Tests für einzelne Schichten kommen dabei weitgehend zu kurz. Diese Beobachtung bezieht sich ausschließlich auf mein derzeitiges Arbeitsumfeld und lässt keine Rückschlüsse auf allgemeine Entwicklungen in der Softwareentwicklung zu.
Darüber hinaus wurde in fachlichen Diskussionen eine von Spotify genutzte Teststrategie~\cite{Spotify18} erörtert, die als Honeycomb-Modell bekannt ist. Diese unterschiedlichen Erfahrungen und Beobachtungen bilden den Ausgangspunkt dieser Arbeit. Es soll untersucht werden, welche Vor- und Nachteile die Testpyramide und das Honeycomb-Modell im Kontext einer Webanwendung bieten.


\subsection{Ziel der Arbeit}
Diese Arbeit vergleicht die Effektivität der Testpyramide und des Testing-Honeycomb-Modells. Dazu wird eine Webanwendung entwickelt, in der mit Test-Driven Development implementiert wurde. Die Anwendung dient als Plattform für die Organisation und Kommunikation zwischen Veranstaltern und Schaustellern. Sie umfasst Module wie das Bewerbungsmanagement, das Vertragsmanagement und das Veranstaltungsmanagement. Veranstalter können Veranstaltungen erstellen, während Schausteller sich bewerben und Verträge abschließen.  Zudem wurde Spring Security integriert, um ein Sicherheitskonzept zu realisieren. Es wird evaluiert, welcher Ansatz hinsichtlich Testqualität, Laufzeit und Komplexit\"at Vorteile bietet. Zur Messung der Testabdeckung werden JaCoCo und PIT eingesetzt. Weiterhin wird untersucht, wie sich der Einsatz einer H2-Datenbank im Vergleich zu produktionsnahen Datenbanken mit Testcontainers auswirkt. Die Ergebnisse dienen als Grundlage f\"ur praxisnahe Empfehlungen.

\subsection{Aufbau der Arbeit}
Die Arbeit ist in mehrere Kapitel gegliedert. Zunächst werden die theoretischen Grundlagen zu Teststrategien, Datenbanktechnologien und Evaluierungstools dargelegt. Im Anschluss wird die Methodik vorgestellt, welche den generellen Ablauf der Untersuchung beschreibt. 

Im Hauptteil der Arbeit wird schrittweise anhand einzelner Features der Webanwendung vorgegangen. Für jedes Feature werden zunächst beide Teststrategien implementiert und deren Laufzeiten gemessen. Die gewonnenen Erfahrungen und Beobachtungen werden dabei dokumentiert und analysiert. Anschließend werden sicherheitsrelevante Aspekte untersucht, bevor die Frontend-Seiten implementiert werden. Danach erfolgt die Einführung von Spring Security, wobei zusätzliche Tests ergänzt werden, falls bestimmte Funktionen zuvor nicht berücksichtigt wurden und nochmal gemessen.

Den Abschluss bildet eine Zusammenfassung der experimentellen Ergebnisse sowie ein Ausblick auf zukünftige Entwicklungen im Bereich der Softwaretests.
